\section{Related Work}
\label{sec:related}

\paragraph{Embedding benchmarks.}
SentEval evaluates sentence representations on classification, natural language inference and sentence similarity tasks \citep{conneau2018senteval}.
USEB, introduced with TSDAE, evaluates unsupervised sentence embeddings on four domain-specific datasets \citep{wang2021tsdae}.
BEIR evaluates retrieval systems zero-shot on 18 datasets \citep{thakur2021beir}.
Our corpus and relevance judgements come from its Natural Questions subset.
MTEB covers 8 tasks and 58 datasets, one of which is retrieval \citep{muennighoff2022mteb}.
These benchmarks score embeddings or ranked passages directly.
None of them includes a downstream reader, and none reports how a retrieval gain changes answer accuracy.

\paragraph{Embedding selection and deployment cost.}
\citet{cirillo2026ge2} compare a hosted embedding model (Google Embeddings~2) with five open-source models on four BEIR subsets and measure per-query latency on CPU.
They report that the hosted model reaches an average nDCG@10 of 0.638 at 231.6\,ms median latency, ``roughly 14$\times$ slower than the fastest local models''.
Their quality measures are retrieval metrics.
\citet{baidya2026choosing} proposes a decision framework for embedding selection by task, latency, cost and deployment constraints.
It reports that ``training objective, not model size, is the dominant factor in retrieval quality''.
It states that it ``does not measure end-to-end RAG answer quality''.

\paragraph{RAG evaluation.}
KILT evaluates knowledge-intensive tasks on downstream performance and on the provenance of the evidence a model uses \citep{petroni2021kilt}.
RAGAS \citep{es2023ragas} defines reference-free metrics for faithfulness, answer relevance and context relevance.
Its evaluation does not compare retrievers.
ARES \citep{saadfalcon2024ares} trains lightweight judges and calibrates them with prediction-powered inference on a small human-annotated set.
It scores RAG systems built from three retrievers and three generators, and each system retrieves one passage.
ARES does not isolate the contribution of retrieval to answer quality and does not measure computational cost.
eRAG \citep{salemi2024erag} evaluates each retrieved document by the reader's downstream output on it.
Its authors report that retrieval evaluation based on query--document relevance labels ``shows a small correlation with the RAG system's downstream performance''.

\paragraph{Retrieval quality and answer quality.}
\citet{leto2024toward} vary the accuracy of approximate nearest-neighbour search in a RAG pipeline.
They report that ``lowering search accuracy has minor implications for RAG performance while potentially increasing retrieval speed and memory efficiency''.
They also report that adding non-gold documents lowers answer correctness relative to a gold-only context.
Their study uses a single dense retriever.
\citet{liu2024lost} show that readers use relevant information less well when it appears in the middle of a long context.
BERGEN \citep{rau2024bergen} benchmarks retrievers, rerankers and LLMs for question answering and compares three retrievers across several LLMs.
RAGPerf \citep{li2026ragperf} profiles RAG systems end to end, collecting throughput, memory and utilisation together with accuracy metrics.
These studies show that retrieval choices affect answers and cost.
They do not report a transfer ratio from a retrieval metric to answer accuracy with confidence intervals.

\paragraph{Positioning.}
We are not aware of work that measures retrieval-to-answer transfer ratios jointly with measured cost on free hardware.
Our study holds one reader, one corpus and one query set fixed and varies seven retrieval methods.
It reports paired tests for every pairwise comparison and measures cost on one free Tesla T4.
